% Relatório da Atividade 5 de Computação Gráfica
% Compilação no Overleaf: pdfLaTeX.
% Envie os PNGs ao lado de main.tex ou na pasta capturas/.
\documentclass[12pt,a4paper]{article}

\usepackage[utf8]{inputenc}
\usepackage[T1]{fontenc}
\usepackage[brazil]{babel}
\usepackage{lmodern}
\usepackage{amsmath}
\usepackage{array}
\usepackage{booktabs}
\usepackage{float}
\usepackage{geometry}
\usepackage{graphicx}
\usepackage{tabularx}
\usepackage[hidelinks]{hyperref}
\geometry{left=3cm,right=2cm,top=2.5cm,bottom=2.5cm}

\newcommand{\printouespaco}[1]{%
  \IfFileExists{#1}{%
    \includegraphics[width=\linewidth]{#1}%
  }{%
    \IfFileExists{capturas/#1}{%
      \includegraphics[width=\linewidth]{capturas/#1}%
    }{%
      \fbox{\parbox[c][4cm][c]{0.90\linewidth}{\centering
        \textbf{INSERIR PRINT AQUI}\\[0.2cm]
        \texttt{#1}}}%
    }%
  }%
}

\newcommand{\captura}[3][0.467\textwidth]{%
  \begin{minipage}[t]{#1}
    \centering
    \printouespaco{#2}%
    \par\smallskip\small #3\par
  \end{minipage}%
}

\begin{document}
\hypersetup{pageanchor=false}
\begin{titlepage}
\centering
{\Large Universidade Federal do ABC\par}
\vspace{1.0cm}
{\large Computação Gráfica\par}
\vspace{3.0cm}
{\LARGE\bfseries Trabalho prático: jogo 3D\par}
\vspace{0.45cm}
{\Large Nave contra asteroides\par}
\vfill
\begin{tabular}{ll}
Cássia dos Santos Coelho & RA: 11202231970\\
Lucas Reis de Oliveira da Silva & RA: 11202230836\\
\end{tabular}
\vfill
{\large \today\par}
\end{titlepage}
\hypersetup{pageanchor=true}

\begin{abstract}
Este trabalho estende o jogo 2D de nave e asteroides para uma cena tridimensional. A nave se move na vertical e na lateral para alinhar disparos com asteroides que se aproximam. A cena inclui duas câmeras em perspectiva, dois tipos de luz e dois objetos texturizados. O programa foi desenvolvido em Python com GLFW, ModernGL, NumPy, PyGLM e Pillow, com shaders GLSL próprios em OpenGL 4.0 Core.
\end{abstract}

\section{Objetivo e funcionamento}

O jogador precisa somar dez acertos antes de perder quatro vidas. Cada acerto em um asteroide vale um ponto. A colisão de um asteroide com a nave ou a sua saída pelo lado esquerdo custa uma vida. Os disparos avançam no eixo $X$, enquanto a nave se desloca em $Y$ e $Z$. Um guia tracejado indica a direção dos disparos. O farol espacial é fixo e serve de referência visual na lateral do setor de jogo.

As dez marcas de pontuação e os quatro quadrados de vida são desenhados com geometria acima da cena. O título da janela informa o estado da partida, a câmera, a iluminação, o VSync e a taxa de quadros.

\section{Arquitetura e organização do código}

O arquivo \texttt{main.py} cria a janela e o contexto gráfico. A cada quadro, ele processa a entrada, chama \texttt{Jogo.atualizar} e usa o estado resultante para desenhar. A simulação mantém posições e velocidades em três dimensões e funciona independentemente do contexto gráfico.

\begin{table}[H]
\centering\small
\begin{tabularx}{\textwidth}{>{\raggedright\arraybackslash}p{5.1cm}X}
\toprule
Arquivo & Responsabilidade \\
\midrule
\texttt{main.py} & Janela, eventos, laço, fundo e HUD. \\
\texttt{Jogo.py} e \texttt{colisoes.py} & Estados, pontuação, vidas e testes contínuos de contato. \\
\texttt{Nave.py} e \texttt{Disparo.py} & Controles de movimento e projéteis. \\
\texttt{Asteroide.py} e \texttt{GerenciaAsteroides.py} & Movimento, pulsação, criação e reposição dos asteroides. \\
\texttt{malha.py} e \texttt{camera.py} & Leitura OFF, normais, enquadramento e câmera. \\
\texttt{cena3d.py} & Geometrias 3D, materiais, texturas e desenho da cena. \\
\texttt{shaders/} & Transformações, cores da camada 2D e iluminação da cena 3D. \\
\bottomrule
\end{tabularx}
\caption{Organização dos módulos.}
\end{table}

O jogo 2D forneceu o laço, os estados, as classes da simulação, o HUD, as geometrias planas e os shaders básicos. Na adaptação 3D, o deslocamento e as colisões ganharam uma terceira coordenada e o shader 2D passou a receber uma matriz de modelo. Do visualizador da apostila foram reaproveitados \texttt{malha.py}, a base de \texttt{camera.py}, o cálculo da matriz normal, o envio de uniformes e os callbacks de redimensionamento e zoom. Os shaders de visualização de normais do protótipo deram lugar aos shaders de iluminação e texturas. O farol, a nave OFF usada na entrega, a malha irregular dos asteroides e os dois mapas de cor foram produzidos para esta etapa.

\section{Laço de jogo e animação}

A cada quadro o programa consulta os eventos, lê \texttt{glfw.get\_time()}, atualiza a simulação, desenha e troca os buffers. O intervalo utilizado é
\begin{equation}
  \Delta t = \min\!\left(\max(t_n-t_{n-1},0),\,0{,}05\right)\ \mathrm{s}.
\end{equation}
O limite reduz saltos quando a janela é arrastada ou o processo fica momentaneamente parado. Os deslocamentos seguem $\Delta\mathbf{p}=\mathbf{v}\Delta t$. A nave usa $1{,}1$ unidade por segundo por eixo, os disparos avançam a $2{,}4$ e os asteroides avançam entre $0{,}30$ e $0{,}43$, com pequena deriva lateral que se inverte nos limites do setor. A fase de cada asteroide aumenta em $2{,}4\Delta t$ e controla a pulsação do raio visual e a rotação da malha.

A tecla V solicita intervalo 1 ou 0 por meio de \texttt{glfw.swap\_interval}. O título é atualizado aproximadamente a cada meio segundo com a razão entre o número de quadros e o tempo real transcorrido.
Em execuções de $800\times800$ com uma GeForce GTX 1060 e GLFW sobre XWayland, a mediana de cinco leituras do título foi de $60{,}0$ FPS com VSync ligado e $2\,251{,}6$ FPS com VSync desligado. A velocidade lateral medida foi de $1{,}101$ e $1{,}100$ unidade por segundo, respectivamente. Na sessão com GLFW Wayland nativo, as medianas foram $1\,699{,}5$ e $1\,641{,}1$ FPS, sem limitação a $60$ FPS. A taxa depende da apresentação de quadros na plataforma, enquanto o deslocamento segue o tempo medido.

Durante a pausa, posições, fases, disparos e enquadramento permanecem fixos. O relógio do laço continua sendo atualizado e o quadro da retomada usa $\Delta t=0$. O mesmo descarte do intervalo vale ao iniciar ou reiniciar a partida.

\section{Colisões e prevenção de tunelamento}

Para o par disparo--asteroide, as posições inicial e final do centro do asteroide são subtraídas das posições correspondentes do disparo. O resultado é o segmento de movimento relativo $\mathbf{r}(u)=\mathbf{r}_0+u(\mathbf{r}_1-\mathbf{r}_0)$. O primeiro contato é a primeira raiz de
\begin{equation}
  \left\|\mathbf{r}(u)\right\|^2
  = \left(r_{\mathrm{disparo}} + r_{\mathrm{asteroide}}\right)^2,
  \qquad 0\leq u\leq 1.
\end{equation}
Quando o par já começa em contato, usa-se $u=0$. Para cada disparo com mais de um alvo possível, o menor $u$ determina qual alvo foi atingido. O disparo é consumido, o asteroide é reposto à frente e a pontuação aumenta.

A nave é aproximada por uma caixa alinhada aos eixos. Para a colisão com o asteroide, a caixa é expandida pelo raio da esfera e o segmento relativo é testado contra ela. O método das \emph{slabs} calcula a entrada e a saída de cada intervalo nos três eixos. A caixa expandida serve de aproximação conservadora, especialmente nos cantos. Usar os trajetos dos centros em ambos os testes permite detectar contatos que ocorrem entre os estados de quadros consecutivos. A colisão com a nave ou a passagem de um asteroide pela borda custa uma vida e provoca a sua reposição.

\section{Pipeline gráfico, transformações e câmera}

A aplicação solicita OpenGL 4.0 Core e os shaders declaram \texttt{\#version 400 core}. A cena 3D reúne nave, asteroides e farol com geometrias diferentes e utiliza teste de profundidade e descarte de faces traseiras. A orientação frontal é anti-horária. VBOs, EBO do HUD, VAOs e texturas são criados na inicialização. Desenhos sucessivos variam matrizes e outros uniformes. A cena utiliza \texttt{TRIANGLES}, \texttt{LINES}, \texttt{LINE\_LOOP} e \texttt{POINTS}. O HUD ocupa um viewport quadrado centralizado, com profundidade desligada, e trabalha diretamente no plano normalizado da tela.

A nave é lida de \texttt{off/nave.off}. Se $\mathbf{c}$ for o centro da caixa envolvente e $r$ a distância ao vértice mais afastado, a matriz $S(0{,}26/r)T(-\mathbf{c})$ centraliza e dimensiona a malha. A translação da nave é aplicada depois. Os dados de posição enviados ao VBO mantêm as coordenadas do OFF. O modelo dos asteroides é composto por translação, rotação e escala em função do raio visual. O farol reutiliza um cubo com translação e escalas diferentes em cada eixo.

O vertex shader calcula $\mathbf{p}_{\mathrm{clip}}=PVM\mathbf{p}_{\mathrm{local}}$. A matriz normal $N_M=(M_{3\times3}^{-1})^T$ é enviada como uniforme separado. A inversa transposta é a transformação que preserva a perpendicularidade das normais após uma escala não uniforme, como as aplicadas às partes do farol.

\begin{table}[H]
\centering\small
\begin{tabularx}{\textwidth}{>{\raggedright\arraybackslash}p{3.0cm}X}
\toprule
Modo & Construção da câmera \\
\midrule
Geral & Alvo fixo em $(0{,}15;-0{,}08;0)$ e olho deslocado na direção normalizada de $(-0{,}23;0{,}58;0{,}91)$, a $4{,}30$ unidades no início. \\
Seguidora & Alvo a $0{,}70$ unidade à frente da nave em $X$. O olho segue a direção normalizada de $(-1;0{,}38;0)$, inicialmente a $2{,}00$ unidades do alvo. \\
\bottomrule
\end{tabularx}
\caption{Câmeras disponíveis durante a partida.}
\end{table}

Nos dois modos, \texttt{glm.lookAt} recebe olho, alvo e vetor para cima $(0;1;0)$. A projeção usa perspectiva, abertura vertical de $45^\circ$, \emph{near} de $0{,}08$ e \emph{far} de $12$. A abertura aumenta em janelas estreitas para preservar o campo horizontal. O plano próximo deixa espaço para a nave na visão seguidora com zoom. O distante cobre os objetos desenhados mesmo com a câmera afastada até $6$ unidades. O callback \texttt{glfw.set\_framebuffer\_size\_callback} ajusta o viewport da cena e a razão largura/altura usada pelas duas projeções.

\section{Iluminação e materiais}

O fragment shader avalia Blinn--Phong com a normal interpolada e renormalizada. Posição do fragmento, olho, vetores de luz e normal estão no espaço do mundo. A parcela ambiente usa a intensidade $0{,}52$ multiplicada por $\kappa_a$. Para cada fonte ativa, a parcela difusa usa $\max(\mathbf{N}\cdot\mathbf{L},0)$ e a especular usa $\max(\mathbf{N}\cdot\mathbf{H},0)^\alpha$ quando a face recebe luz, com $\mathbf{H}$ sendo o vetor médio normalizado entre a direção da luz e a direção do observador.

A direção fixa da primeira fonte é proporcional a $(0{,}35;0{,}70;0{,}62)$ e a cor RGB linear é $(0{,}94;0{,}94;1)$. A luz pontual acompanha a nave com deslocamento $(0{,}32;0{,}40;0{,}18)$ e cor $(0{,}55;0{,}88;1)$. A atenuação pontual é $1/(1+0{,}23d+0{,}20d^2)$, em que $d$ é a distância até a fonte. A tecla L permite observar as duas juntas e cada uma separadamente.

\begin{table}[H]
\centering\small
\begin{tabular}{lrrrr}
\toprule
Material & $\kappa_a$ & $\kappa_d$ & $\kappa_s$ & $\alpha$ \\
\midrule
Metal da nave & 0,25 & 0,76 & 0,53 & 44 \\
Rocha do asteroide & 0,29 & 0,95 & 0,10 & 12 \\
Painel do farol & 0,28 & 0,81 & 0,27 & 28 \\
Sinal do farol & 0,30 & 0,46 & 0,26 & 26 \\
\bottomrule
\end{tabular}
\caption{Coeficientes dos materiais principais. O sinal também tem emissão própria.}
\end{table}

\section{Texturas e procedência dos recursos}

\texttt{rocha.png} é um mapa de cor de $256\times256$ pixels, produzido com ruído de várias escalas e marcas de crateras. As UVs são geradas com a malha, usando longitude e latitude. As normais dos asteroides são obtidas com o cálculo por vértice reaproveitado do visualizador. A costura repete posição e normal em $u=0$ e $u=1$, e as coordenadas horizontais dos polos são definidas por setor. A rocha usa \texttt{GL\_REPEAT} em $u$ e \texttt{GL\_CLAMP\_TO\_EDGE} em $v$.

\texttt{paineis.png} é um mapa de cor de $128\times128$ pixels desenhado para o farol. O mapeamento é planar e percorre o intervalo de $0$ a $2$ no eixo $v$ de cada face. A textura repete nos dois eixos. Ambos os PNGs são carregados com Pillow, em RGBA, com inversão vertical. O formato interno \texttt{GL\_SRGB8\_ALPHA8} converte a amostra para RGB linear antes da iluminação. O fragment shader aplica a curva sRGB uma vez na saída 3D. A minificação é trilinear, com mipmaps gerados na inicialização, e a magnificação é linear.

A nave OFF, as geometrias procedurais e os dois PNGs são recursos originais desta versão, produzidos com apoio do Codex para o trabalho. A fonte é o próprio projeto. Os recursos originais acompanham esta entrega acadêmica, sem licença de terceiros aplicável. O protótipo de Cássia havia usado o modelo \href{https://sketchfab.com/3d-models/spaceship-1efc26312afe4d4c907ffcdba12f4a6a}{\emph{Spaceship}, de madalenacardoso}, convertido de GLB para OFF com MeshLab, e uma esfera de origem ainda não identificada. Esses arquivos foram substituídos na versão entregue. O código de partida está registrado no \href{https://github.com/cassia-santos/Atividade_5_CG/tree/26ae2dadadea36136a7d075f83f99ab339ec5f45/prototipo1}{repositório do protótipo}.

\section{Estados e controles}

\begin{table}[H]
\centering\small
\begin{tabularx}{\textwidth}{>{\raggedright\arraybackslash}p{3.2cm}X}
\toprule
Estado & Comportamento \\
\midrule
Início & Exibe o símbolo de reprodução. Enter começa a partida. \\
Andamento & Atualiza a nave, os disparos, os asteroides e as colisões. \\
Pausa & Congela a simulação e a câmera. P retoma. \\
Fim & Exibe o símbolo verde de vitória ou o X de derrota. Enter reinicia. \\
\bottomrule
\end{tabularx}
\caption{Estados da partida.}
\end{table}

As setas para cima e para baixo movem a nave verticalmente, e as setas para a esquerda e para a direita movem na lateral. O movimento contínuo usa \texttt{glfw.get\_key}. Enter inicia ou reinicia, Espaço dispara, P alterna pausa, C alterna câmeras, L alterna as luzes, V alterna VSync e Esc fecha a janela. Essas ações usam o callback de teclado. A roda do mouse ajusta a distância pelo callback de rolagem.

\clearpage
\section{Capturas de execução}

\begin{figure}[H]
\centering
\captura{print-inicio.png}{Início, com o símbolo de reprodução.}\hfill
\captura{print-andamento.png}{Andamento com um ponto, visão geral e ambas as luzes.}
\par\medskip
\captura{print-pausa.png}{Pausa, com duas barras centrais.}\hfill
\captura{print-fim.png}{Fim de partida por perda das quatro vidas.}
\caption{Os quatro estados da partida.}
\end{figure}

\clearpage
\begin{figure}[H]
\centering
\captura[0.315\textwidth]{print-direcional.png}{Luz direcional.}\hfill
\captura[0.315\textwidth]{print-pontual.png}{Luz pontual.}\hfill
\captura[0.315\textwidth]{print-seguidora.png}{Câmera seguidora, ambas as luzes.}
\caption{Câmera seguidora e comparação das fontes. A rocha e os painéis do farol podem ser vistos com iluminação na cena.}
\end{figure}

\section{Divisão de tarefas}

\textbf{Cássia dos Santos Coelho.} Preparou o protótipo 3D publicado no repositório. Adaptou as classes da etapa 2D para incluir a profundidade, integrou os modelos OFF e os buffers 3D, manteve o HUD e os shaders 2D, acrescentou as duas câmeras e programou o zoom. Registrou para o relatório a origem e a conversão da nave do protótipo e listou os módulos e funções reaproveitados.

\smallskip
\noindent\textbf{Lucas Reis de Oliveira da Silva.} Conduziu a etapa de conclusão, reuniu o enunciado e o material de partida, definiu o escopo de complementação do protótipo e solicitou a implementação, a validação e o relatório com apoio do Codex. Forneceu o relatório anterior como referência para a estrutura do texto e da capa.

\section{Declaração de uso de inteligência artificial}

Foi utilizado o OpenAI Codex para comparar o protótipo com o enunciado, complementar geometrias, shaders, iluminação, texturas e integração da cena, executar verificações, obter as capturas e redigir o relatório em LaTeX. A responsabilidade técnica e intelectual sobre o código e o relatório submetidos é da dupla.

\section{Conclusão}

O protótipo passou a exibir três tipos de geometria em perspectiva, com modelos iluminados e duas texturas. A etapa preservou as regras da partida, o laço por tempo decorrido e os testes contínuos de contato. As duas câmeras e a seleção de luzes permitem observar a mesma cena sob enquadramentos e condições de iluminação diferentes.

\end{document}
